% GENERATED FILE. Edit canonical lessons, then run npm run generate:books.
% canonical-source-hash: fnv1a64:c3dfa75bde42c0f3
% canonical-lessons: SA-C215-calah, SA-C215-mandah, SA-C215-tiksnah, SA-C215-kunthah, SA-C215-mulyavan

\chapter{Moving, Slow, Sharp, Blunt, Valuable}
\label{ch:sa-moving-slow-sharp-blunt-valuable}
\hlchaptermodality{215}

\begin{chapteropening}
\textbf{By the end of this chapter:} \emph{I can say moving, slow, sharp, blunt and valuable.}

Complete the last lesson of chapter 215: I can say moving, slow, sharp, blunt and valuable.
\end{chapteropening}

\section[\texorpdfstring{\sk{चलः} (calaḥ)}{चलः (calaḥ)}]{\sk{चलः} (calaḥ) — moving, unsteady}

\label{lesson:SA-C215-calah}



\begin{quote}
Before the new one: say the Sanskrit for without, then the Sanskrit for towards.
\end{quote}

\subsection*{You'll want to know: \sk{चलः}}
\textbf{\sk{चलः}} — \emph{calaḥ} — \textquotedblleft{}moving, unsteady\textquotedblright{}.

\begin{grammarlens}[title={What you've built}]
One more word for this chapter.
\end{grammarlens}

\subsection*{Your turn}
\begin{itemize}
\raggedright
  \item \emph{Say it:} \emph{calaḥ}
  \item \emph{Say it:} \emph{calaḥ}, once more
  \item \emph{Say it:} say \emph{calaḥ}
  \item \emph{Recall:} say the Sanskrit for without, then the Sanskrit for towards, then say \emph{calaḥ} again
\end{itemize}

\begin{tcolorbox}[breakable,colback=teal!4,colframe=teal!35!black,title={Before you move on}]
What does \sk{चलः} mean? (\textquotedblleft{}Moving, unsteady\textquotedblright{}.) Say it once more.
\end{tcolorbox}
\section[\texorpdfstring{\sk{मन्दः} (mandaḥ)}{मन्दः (mandaḥ)}]{\sk{मन्दः} (mandaḥ) — slow, dull}

\label{lesson:SA-C215-mandah}



\begin{quote}
Before the new one: say the Sanskrit for towards, then the Sanskrit for moving.
\end{quote}

\subsection*{You'll want to know: \sk{मन्दः}}
\textbf{\sk{मन्दः}} — \emph{mandaḥ} — \textquotedblleft{}slow, dull\textquotedblright{}.

\begin{grammarlens}[title={What you've built}]
One more word for this chapter.
\end{grammarlens}

\subsection*{Your turn}
\begin{itemize}
\raggedright
  \item \emph{Say it:} \emph{mandaḥ}
  \item \emph{Say it:} \emph{mandaḥ}, once more
  \item \emph{Say it:} say \emph{mandaḥ}
  \item \emph{Recall:} say the Sanskrit for towards, then the Sanskrit for moving, then say \emph{mandaḥ} again
\end{itemize}

\begin{tcolorbox}[breakable,colback=teal!4,colframe=teal!35!black,title={Before you move on}]
What does \sk{मन्दः} mean? (\textquotedblleft{}Slow, dull\textquotedblright{}.) Say it once more.
\end{tcolorbox}
\section[\texorpdfstring{\sk{तीक्ष्णः} (tīkṣṇaḥ)}{तीक्ष्णः (tīkṣṇaḥ)}]{\sk{तीक्ष्णः} (tīkṣṇaḥ) — sharp}

\label{lesson:SA-C215-tiksnah}



\begin{quote}
Before the new one: say the Sanskrit for moving, then the Sanskrit for slow.
\end{quote}

\subsection*{You'll want to know: \sk{तीक्ष्णः}}
\textbf{\sk{तीक्ष्णः}} — \emph{tīkṣṇaḥ} — \textquotedblleft{}sharp\textquotedblright{}.

\begin{grammarlens}[title={What you've built}]
One more word for this chapter.
\end{grammarlens}

\subsection*{Your turn}
\begin{itemize}
\raggedright
  \item \emph{Say it:} \emph{tīkṣṇaḥ}
  \item \emph{Say it:} \emph{tīkṣṇaḥ}, once more
  \item \emph{Say it:} say \emph{tīkṣṇaḥ}
  \item \emph{Recall:} say the Sanskrit for moving, then the Sanskrit for slow, then say \emph{tīkṣṇaḥ} again
\end{itemize}

\begin{tcolorbox}[breakable,colback=teal!4,colframe=teal!35!black,title={Before you move on}]
What does \sk{तीक्ष्णः} mean? (\textquotedblleft{}Sharp\textquotedblright{}.) Say it once more.
\end{tcolorbox}
\section[\texorpdfstring{\sk{कुण्ठः} (kuṇṭhaḥ)}{कुण्ठः (kuṇṭhaḥ)}]{\sk{कुण्ठः} (kuṇṭhaḥ) — blunt}

\label{lesson:SA-C215-kunthah}



\begin{quote}
Before the new one: say the Sanskrit for slow, then the Sanskrit for sharp.
\end{quote}

\subsection*{You'll want to know: \sk{कुण्ठः}}
\textbf{\sk{कुण्ठः}} — \emph{kuṇṭhaḥ} — \textquotedblleft{}blunt\textquotedblright{}.

\begin{grammarlens}[title={What you've built}]
One more word for this chapter.
\end{grammarlens}

\subsection*{Your turn}
\begin{itemize}
\raggedright
  \item \emph{Say it:} \emph{kuṇṭhaḥ}
  \item \emph{Say it:} \emph{kuṇṭhaḥ}, once more
  \item \emph{Say it:} say \emph{kuṇṭhaḥ}
  \item \emph{Recall:} say the Sanskrit for slow, then the Sanskrit for sharp, then say \emph{kuṇṭhaḥ} again
\end{itemize}

\begin{tcolorbox}[breakable,colback=teal!4,colframe=teal!35!black,title={Before you move on}]
What does \sk{कुण्ठः} mean? (\textquotedblleft{}Blunt\textquotedblright{}.) Say it once more.
\end{tcolorbox}
\section[\texorpdfstring{\sk{मूल्यवान्} (mūlyavān)}{मूल्यवान् (mūlyavān)}]{\sk{मूल्यवान्} (mūlyavān) — valuable}

\label{lesson:SA-C215-mulyavan}



\begin{quote}
Before the new one: say the Sanskrit for sharp, then the Sanskrit for blunt.
\end{quote}

\subsection*{You'll want to know: \sk{मूल्यवान्}}
\textbf{\sk{मूल्यवान्}} — \emph{mūlyavān} — \textquotedblleft{}valuable\textquotedblright{}.

\begin{grammarlens}[title={What you've built}]
One more word for this chapter.
\end{grammarlens}

\subsection*{Your turn}
\begin{itemize}
\raggedright
  \item \emph{Say it:} \emph{mūlyavān}
  \item \emph{Say it:} \emph{mūlyavān}, once more
  \item \emph{Say it:} say \emph{mūlyavān}
  \item \emph{Recall:} say the Sanskrit for sharp, then the Sanskrit for blunt, then say \emph{mūlyavān} again
\end{itemize}

\begin{tcolorbox}[breakable,colback=teal!4,colframe=teal!35!black,title={Before you move on}]
What does \sk{मूल्यवान्} mean? (\textquotedblleft{}Valuable\textquotedblright{}.) Say it once more.
\end{tcolorbox}
